\chapter{因果路径与安全投影}
\label{ch:safe}

同一次蒸汽指令可以同时加热冷点、打湿表面、并在更长的时间上改变结壳热阻。把这三件事合成一个“蒸汽有利”的分数，会把薄件上的冷凝收益写成厚件上的菜单。安全层要解决的是另一件事：品质模型给出的参考动作，在进入四个执行器之前必须被只依赖硬测量的不等式改写。水位、门、剩余电流和加热器温度构成这条投影的自变量；进度估计留在品质回路。IEC 60335-1 与 GB 4706.1 要求这些测量点存在，条款符合性需要独立的型式试验。

\section{结构方程与干预}

决策状态的分量满足结构方程
\begin{equation}
\label{eq:scm}
Z_i=f_i(\mathrm{Pa}_i,U_i),\qquad i=1,\ldots,n_z.
\end{equation}
父节点 \(\mathrm{Pa}_i\) 由机理指定，见表\ref{tab:scm}。\(U_i\) 是外生噪声。学习阶段只调整函数的参数和事先声明的可选边；一次烹饪的短序列上，能量守恒的边保持固定。

\begin{table}[htbp]
\centering
\caption{因果骨架。内部热阻既是参数也是子节点，因为结壳会增大它。}
\label{tab:scm}
\begin{tabularx}{\textwidth}{llX}
\toprule
子节点 & 父节点 & 物理内容 \\
\midrule
界面蒸汽压 & 蒸汽流、排汽、腔温 & 腔体水分平衡 \\
表面温度 & 蒸汽压、腔温、风速、加热、内部温度 & \eqref{eq:twonode} 第一式 \\
内部温度 & 表面温度、内部热阻、内部热容 & \eqref{eq:twonode} 第二式 \\
表面含水 & 蒸汽压、表面温度、传质系数 & 蒸发与冷凝 \\
冷点进度 & 内部温度的历史 & \eqref{eq:F} \\
表面褐变 & 表面温度、表面含水 & \eqref{eq:kinetics} \\
结壳 & 蛋白进度、表面含水 & \eqref{eq:structure} \\
内部热阻 & 结壳、厚度 & 结构到传递 \\
\bottomrule
\end{tabularx}
\end{table}

观测条件分布 \(p(\alpha_m\mid \dot m_{\mathrm{st}})\) 混有控制器对当前低温的反应：没熟所以多喷蒸汽。干预分布 \(p(\alpha_m\mid\mathrm{do}(\dot m_{\mathrm{st}}=m^\star))\) 把决定蒸汽流量的策略从结构方程中删去，只保留蒸汽到界面蒸汽压的物理边\cite{pearl2009}。孪生的积分实现干预：替换执行器轨迹后再积分。反事实还需要先用已经发生的观测反推外生噪声。本次没有反事实抽样，界面因此只显示已经写出的通量代理。S1 与档位研究里的四个策略是四条不同的闭环，策略之间的 \(e_{\alpha}\) 差描述闭环差别，路径偏效应需要在同一噪声下替换蒸汽再积分。

\section{三条路径}

蒸汽对品质的全导数分成冷点项、表面项和慢热阻项：
\begin{align}
\text{路径 A: }& \dot m_{\mathrm{st}}\to p_v\to q_{\mathrm{lat}}\to T_2\to\alpha_m\to Q,\\
\text{路径 B: }& \dot m_{\mathrm{st}}\to p_v\to\varepsilon_{l,s}\to(\alpha_b,\,m_{\mathrm{loss}})\to Q,\\
\text{路径 C: }& \dot m_{\mathrm{st}}\to\varepsilon_{l,s}\to S\to R_{\mathrm{int}}\to T_2\to\alpha_m\to Q.
\end{align}
路径 A 在冷点落后时为正。路径 B 在表面已经湿、且后续需要烤色时为负。路径 C 是慢路径：短探针窗口内增量是 \(O(\tau/\tau_S)\)，全时域必须保留。总效应
\begin{equation}
\label{eq:path}
\frac{\mathrm{d}Q}{\mathrm{d}m_{\mathrm{st}}}
=\underbrace{\frac{\partial Q}{\partial\alpha_m}\frac{\partial\alpha_m}{\partial m_{\mathrm{st}}}\Big|_{\!S}}_{\text{A}}
+\underbrace{\frac{\partial Q}{\partial\varepsilon_{l,s}}\frac{\partial\varepsilon_{l,s}}{\partial m_{\mathrm{st}}}}_{\text{B}}
+\underbrace{\frac{\partial Q}{\partial\alpha_m}\frac{\partial\alpha_m}{\partial R_{\mathrm{int}}}\frac{\partial R_{\mathrm{int}}}{\partial m_{\mathrm{st}}}}_{\text{C}}.
\end{equation}

\begin{proposition}[路径加性]
若品质泛函 \(Q\) 只通过 \(\alpha_m\)、\(\varepsilon_{l,s}\) 和 \(R_{\mathrm{int}}\) 依赖蒸汽质量，则全导数恰好分解为\eqref{eq:path} 的三项。
\end{proposition}

\begin{proof}
链式法则在所述依赖下把全导数写成三条互不重叠的通路。其余状态若只作为这三者的中间变量，已被吸收进偏导数。
\end{proof}

\begin{figure}[htbp]
\centering
\includegraphics[width=0.88\textwidth]{path_abc.png}
\caption{同一次蒸汽指令的三条路径。实线对应本次通量里有对应项的边；路径 C 经过结壳，结壳不改十六条控制体的导热，因而在积分里没有闭合。}
\label{fig:path}
\end{figure}

图\ref{fig:mpaths}把同一分叉画成界面上的三条去向。路径 A 把潜热经表面传到冷点，路径 B 改表面含水，路径 C 只进入结壳参照。规格里的 C 还要降低导热，核没有这一步。

\begin{figure}[htbp]
\centering
\includegraphics[width=0.86\textwidth]{M_paths.png}
\caption{蒸汽离开界面之后的三条去向。路径 C 在核里不乘回导热。}
\label{fig:mpaths}
\end{figure}

图\ref{fig:path}给出分叉。本次积分闭合了潜热通量对蒸汽、活度和风机的依赖，以及表面失水对活度和风机的依赖。结壳只累加，不乘进导热，所以路径 C 在多数 S1 样本上没有数值。界面通量对蒸汽指令的偏导数
\begin{equation}
\label{eq:dqlat}
\frac{\partial q_{\mathrm{lat}}}{\partial u_{\mathrm{st}}}
=18000\,(1-a_w)+9000\max(a_w-0.35,0)\,u_v
\end{equation}
在种子 7、活度被裁到 0.99、风机 0.4 时等于 \(2484\,\mathrm{W}/\mathrm{m}^2\)。它为正，主要来自关掉蒸汽时的蒸发冷却被撤掉：饱和含水上，蒸汽的热效应是撤冷却。种子 7 探针窗内冷点从 \(8.01^\circ\mathrm{C}\) 到 \(8.05^\circ\mathrm{C}\)，短窗内 A 与 C 的温度增量很小。浏览器侧栏的三根条来自该帧已经写出的通量代理，\eqref{eq:path} 的偏效应积分留在规格里。

厚度增大使 A 变小。结壳使 B 更负，并使 C 在烤制段为负。解释输出应当是这三项的数值与区间。厚度、初温和开门同时影响蒸汽用量和最终品质，必须出现在结构里；调味转化和微生物初始数量的绝对值不进入解释语句。

\section{品质参考与安全集}

规格中的有限时域代价同时惩罚未达到的热力进度、表面偏离、不均匀和能量\cite{rawlings2017}。本次加热律是第\ref{ch:ati}章的饱和线性请求，只看见估计温度、方差和厚度。能量被记录为结果列。无探针和对偶两支共用这条加热律，有限时域求解器留在规格里。

品质层的输出记为 \(\vect{u}_{\mathrm{AI}}\)。安全层求解
\begin{equation}
\label{eq:cbf}
\begin{aligned}
\vect{u}={}&\argmin_{\vect{u}}\;\|\vect{u}-\vect{u}_{\mathrm{AI}}\|_W^2\\
\mathrm{s.t.}\quad
&\dot h_j+\gamma_j h_j\ge 0.
\end{aligned}
\end{equation}
\(h_j\ge 0\) 表示安全集内部\cite{ames2017cbf}。自变量只含漏电、加热器温度、水位或空烧温升、腔压、门状态和执行器电流，不含 \(\alpha_m\)。屏障在估计器崩溃时仍然成立。漏电与干烧另有独立断电回路，不读取品质后验，也不经过产生参考动作的处理器。投影负责提前限制，断电负责投影失效后的切断。家用电器对漏电、温升和异常工作的通用要求见 IEC 60335-1 与 GB 4706.1\cite{iec60335,gb4706}。水位、门、剩余电流和加热器温度因此进入安全投影，进度估计留在品质回路。条款符合性需要独立的型式试验。

本次落实的削减是\eqref{eq:hard}。加热器满功率平衡在 \(196.5^\circ\mathrm{C}\)，战役里 \(230^\circ\mathrm{C}\) 削减保持待命。断言另用 \(210^\circ\mathrm{C}\) 的测试函数把请求剪到 \(0.25\)，并且删除后验参数，两个不同后验得到同一个剪裁。两条规则都与后验无关；两条规则在本次加热器轨迹上都停在平衡点之下。要在主表里看见屏障改写指令，需要另设初始温度或增益。图\ref{fig:heater7}给出种子 7 的加热器温度与腔温，平衡点清晰可见。

\begin{figure}[htbp]
\centering
\includegraphics[width=0.84\textwidth]{P_heater.png}
\caption{种子 7 加热器温度与腔温。满功率段停在 \(196.5^\circ\mathrm{C}\)，低于 \(230^\circ\mathrm{C}\) 阈值。}
\label{fig:heater7}
\end{figure}

置信度低于阈值时，送入投影的参考改为预先存储的低功率保持，学习模块的最后一次输出不再进入执行器。核里的保持请求是 \(0.15\)，仍然经过同一条加热器更新。种子 7 在 1872 s 切入保持，加热器从平衡点回落到约 \(90^\circ\mathrm{C}\)，剖面在此后被拉平。评价用停机时刻的 \(\alpha_m\)：终点进度已经饱和到 1，那是保持段的结果。

品质回路读表面温度和平均含水，写出四个请求；安全回路读水位、门、剩余电流和加热器温度，在请求与指令之间改写这四个量。投影可以减小加热，断电可以把四个量同时写成零。种子 7 的工作点上加热器停在 \(196.5^\circ\mathrm{C}\)，请求与指令重合，投影保持待命。要把投影看见，需要另设加热器初温或增益，让管温越过 \(230^\circ\mathrm{C}\)。断言用 \(210^\circ\mathrm{C}\) 的测试函数已经核对：两个不同后验得到同一个剪裁，可行域与后验不相交。IEC 60335-1 与 GB 4706.1 要求这些测量点存在；正文给出的是测量隔离的结构，条款符合性需要独立的型式试验。第\ref{ch:demo}章把请求与指令分两列画在侧栏。

解释句写成“蒸汽有利”或“蒸汽有害”之前，先给出路径 A、B、C 的符号和区间。路径 A 把冷凝潜热经表面传到冷点，短窗口内增量很小，种子 7 在 40 s 里冷点只升 \(0.04\,\mathrm{K}\)；路径 B 改表面含水，探针窗内蒸汽为 1 时蒸发冷却被撤掉；路径 C 经过结壳，核里结壳只累加、不乘回导热，多数 S1 样本上接近零。厚度增大使 A 变小，结壳使 B 更负。本次日志给出的是该帧通量代理，偏效应积分留在规格里。

\FloatBarrier
